\section{Attribution dữ liệu}
\label{sec:ch07-attribution-du-lieu}

\index{attribution dữ liệu}%
\index{hàm ảnh hưởng}%
Lưới ở mục trước đặt ba kỹ thuật cạnh nhánh thứ hai; mục này kiểm lại từng ô của
cột ấy.
\tn{Attribution dữ liệu}{data attribution} hỏi mỗi điểm trong tập huấn luyện đã
đẩy đầu ra tại $x$ đi bao xa, nên nó nhìn vào giai đoạn huấn luyện chứ không
nhìn vào một lượt suy diễn~\autocite{p18unified}. Điểm số $\psi_i(x)$ vì thế
không đo đầu vào nào quan trọng, mà đo dữ liệu nào đã tạo ra hành vi đang xét.

Dạng phép nhiễu ở đây là leave-one-out: bỏ điểm huấn luyện thứ $i$, huấn luyện
lại mô hình, rồi lấy chênh lệch hiệu năng làm điểm số. Ý tưởng có gốc trong
thống kê từ lâu trước khi có mạng sâu, và nó cho một cách đọc phản thực rất
sạch. Cái giá là phải huấn luyện lại một lần cho mỗi điểm dữ
liệu~\autocite{p18unified}. Ngoài ra, bỏ từng điểm một sẽ bỏ sót tương tác: hai
điểm gần trùng nhau có thể cùng đỡ một dự đoán, nên bỏ riêng từng điểm thì đầu
ra hầu như không đổi, còn bỏ cả hai thì đổi hẳn.

Data Shapley xử lý đúng chỗ hở đó bằng cách gọi lại lý thuyết trò chơi hợp tác:
điểm số của một điểm huấn luyện là đóng góp biên trung bình của nó trên mọi tập
con của tập huấn luyện~\autocite{p18unified}. Bài toán tổ hợp cũng lặp lại theo:
với $n$ điểm dữ liệu có $2^n$ tập con, nên các biến thể thực dụng đều là phép
xấp xỉ, cắt bớt số hoán vị được lấy mẫu, thay mô hình gốc bằng một mô hình thay
thế rẻ hơn, hoặc đổi trọng số của các tập con theo kích thước.

Kỹ thuật gradient vào nhánh này qua \tn{hàm ảnh hưởng}{influence function}. Thay
vì huấn luyện lại khi bỏ một điểm, ta hỏi tham số tối ưu dịch chuyển thế nào khi
tăng nhẹ trọng số của điểm đó, rồi khai triển Taylor quanh tham số $\theta^{*}$
đã hội tụ. Ảnh hưởng của điểm huấn luyện $x^{(i)}$ lên mất mát tại đầu vào kiểm
tra $x$ khi ấy viết được thành

\begin{equation}
\label{eq:ch07-ham-anh-huong}
\mathcal{I}(x^{(i)}, x)
= -\,\nabla_{\theta}\mathcal{L}(x,\theta^{*})^{\top}
   H_{\theta^{*}}^{-1}\,
   \nabla_{\theta}\mathcal{L}(x^{(i)},\theta^{*}),
\end{equation}

trong đó $H_{\theta^{*}}$ là ma trận Hessian của rủi ro thực nghiệm tại
$\theta^{*}$; điểm attribution dữ liệu là số đối của đại lượng
này~\autocite{p18unified}.

Phương trình~\ref{eq:ch07-ham-anh-huong} rẻ hơn leave-one-out, và chi phí giảm
đi đó được đánh đổi bằng giả thiết. Khai triển cần hàm mất mát khả vi hai lần và
lồi chặt, hai điều kiện mà học sâu thường không thỏa. Hessian cũng có thể không
xác định dương, nên phải thêm một lượng giảm chấn, và lượng ấy tác động lên độ
chính xác của ước lượng ảnh hưởng~\autocite{p18unified}. Ta gặp lại đúng cấu
trúc của vấn đề mà chương~\ref{ch:attribution-tu-nguyen-ly-dau} đã nêu dưới tên
độ đo: một lựa chọn nằm bên trong phép tính, được biện minh bằng tính khả thi
chứ không bằng miền ứng dụng.

Hai lối đi còn lại nhắm vào cùng hai điểm yếu ấy. Nhóm thứ nhất theo dõi cả
đường huấn luyện thay vì chỉ tham số cuối: TracIn cộng tích vô hướng giữa
gradient của điểm huấn luyện và gradient của điểm kiểm tra qua từng bước cập
nhật, nhờ vậy hai điểm trùng nhau không còn nhận cùng một điểm số, đổi lại phải
lưu các checkpoint trung gian~\autocite{p18unified}. Nhóm thứ hai dùng
\tn{xấp xỉ tuyến tính}{linear approximation}: Datamodel khớp một mô hình tuyến
tính $n$ hệ số nhận vectơ nhị phân
đánh dấu tập con huấn luyện làm đầu vào và dự đoán $f(x)$ làm đầu ra, rồi đọc hệ
số làm điểm số; TRAK ước lượng Datamodel trong một không gian đã biến đổi để bài
toán trở nên lồi~\autocite{p18unified}.

Ba kỹ thuật đã lặp lại đủ, và cùng với chúng là cả ba loại lựa chọn tự do: lấy
mẫu bao nhiêu tập con, xấp xỉ Hessian ra sao, lưu bao nhiêu checkpoint. Nhánh
thứ ba đưa cùng bộ kỹ thuật ấy vào bên trong mô hình.
